前面的凸优化建立了一种强结论：只要问题保持凸性，局部下降、最优性条件和对偶下界就能共同指向全局答案。深度网络、矩阵分解和策略优化不再拥有这种结构，但它们仍大量使用梯度方法。此时更新式可能保持不变，我们却必须调整对算法能够证明什么的期待。

本单元沿六个层次走出凸性：

\begin{enumerate}
\tightlist
\item
  下降引理在非凸地形中只保证接近驻点：复用光滑性结论，推导梯度范数的复杂度界，并说明小梯度不是全局证书。
\item
  二阶结构区分局部极小与鞍点：用 Hessian 的正负曲率解释负曲率逃逸、退化驻点和近似二阶驻点。
\item
  PL 条件让部分非凸问题恢复全局速率：展示凸性并非线性收敛的唯一来源。
\item
  SGD 在期望中下降但噪声会留下平台：把无偏梯度、方差、batch size 与步长写进同一条不等式。
\item
  动量与自适应缩放改变的是动力学：分析曲率方向上的二阶递推，并说明 Adam 的有效坐标步长为何需要额外收敛条件。
\item
  优化误差、统计误差与泛化不能混为一谈：把训练目标下降放回经验风险与总体风险的完整分解。
\end{enumerate}

这里要建立证据等级：目标下降、梯度变小、到达局部极小、找到全局最优、经验风险降低和总体风险改善是六个不同结论。每次报告``收敛''时，都应明确究竟证明或观察到了哪一个。

\subsection{一个简单函数暴露六层差异}

令 \(f(x)=x^4-3x^3+2x^2+1\)。这个四次函数非凸（Hessian 在 \(x<1/3\) 与 \(x>4/3\) 为正，中间为负）。从不同初始点出发：\(x_0=3\) 降落到 \(x\approx2.37\)（局部极小，\(f\approx-0.63\)）；\(x_0=-1\) 降落到 \(x\approx-0.22\)（局部极小，\(f\approx1.04\)）。梯度下降从 3 出发停在局部极小，梯度 < \(10^{-6}\)，但与全局最优（\(x\approx2.37\) 附近）之间仍有 gap。即使在这个光滑一维例子上，梯度消失只保证驻点，不保证全局最优。加入随机梯度噪声后，固定步长会在驻点附近形成大约 \(O(\eta)\) 的平台，梯度不会再降到任意小。

\subsection{怎么读这一章}

离开凸性后，先把承诺降到正确层级。前三篇区分驻点、鞍点和仍可恢复全局速率的特殊结构；后三篇讨论随机梯度、动量与泛化，说明训练曲线变好不等于统计问题已经解决。读时始终分清算法找到了什么、理论证明了什么。凸情形的基准可回看 \hyperref[ch:076]{``凸优化算法：利用梯度、曲率与可分结构''}。

\subsection{本章篇目}

以下篇目按概念依赖排列。

\begin{itemize}
\tightlist
\item \hyperref[sec:09-Convex-Optimization/08-Nonconvex-and-Stochastic-Optimization/01]{``下降引理在非凸地形中只保证接近驻点''}；
\item \hyperref[sec:09-Convex-Optimization/08-Nonconvex-and-Stochastic-Optimization/02]{``二阶结构区分局部极小与鞍点''}；
\item \hyperref[sec:09-Convex-Optimization/08-Nonconvex-and-Stochastic-Optimization/03]{``PL 条件让部分非凸问题恢复全局速率''}；
\item \hyperref[sec:09-Convex-Optimization/08-Nonconvex-and-Stochastic-Optimization/04]{``SGD 在期望中下降但噪声会留下平台''}；
\item \hyperref[sec:09-Convex-Optimization/08-Nonconvex-and-Stochastic-Optimization/05]{``动量与自适应缩放改变的是动力学''}；
\item \hyperref[sec:09-Convex-Optimization/08-Nonconvex-and-Stochastic-Optimization/06]{``优化误差、统计误差与泛化不能混为一谈''}。
\end{itemize}
